\section*{Chapter \ref{ch:g8:ohms-law} --- Resistance and Ohm's Law}

\begin{solution}{exo:g8:ohms-law:1}
How strongly a component opposes the current; \omterm{def:g6:measurement-in-science:measuring}{unit} the \omterm{def:g8:ohms-law:resistance}{ohm},
symbol $\Omega$. $U = R I$; \ $I = U/R$; \ $R = U/I$.
\end{solution}

\begin{solution}{exo:g8:ohms-law:2}
$I = 6 \div 30 = \qty{0.2}{A}$; \ $U = 30 \times 0.5 =
\qty{15}{V}$.
\end{solution}

\begin{solution}{exo:g8:ohms-law:3}
$R = 4.5 \div 0.09 = 50\,\Omega$.
\end{solution}

\begin{solution}{exo:g8:ohms-law:4}
A straight line through the origin --- proportionality. A
steeper line means more \omterm{def:g8:voltage-voltmeter:voltage}{volts} needed per \omterm{def:g8:intensity-ammeter:intensity}{ampere}: greater
\omterm{def:g8:ohms-law:resistance}{resistance}.
\end{solution}

\begin{solution}{exo:g8:ohms-law:5}
$3.0 \div 0.20 = 15$; $6.0 \div 0.40 = 15$. The sameness of the
ratio across rows \emph{is} the law: one constant $R$ serving
every reading is what makes ``the \omterm{def:g8:ohms-law:resistance}{resistance}'' a property of
the component.
\end{solution}

\begin{solution}{exo:g8:ohms-law:6}
The smaller \omterm{def:g8:ohms-law:resistance}{resistance} drinks more: $230 \div 25 =
\qty{9.2}{A}$ against $230 \div 35 \approx \qty{6.6}{A}$.
\end{solution}

\begin{solution}{exo:g8:ohms-law:7}
Series lamps share one $I$; each lamp's \omterm{def:g8:voltage-voltmeter:voltage}{voltage} is $U = R I$
with the same $I$ --- so the larger $R$ multiplies into the
larger $U$: the stronger opposer takes the bigger share.
\end{solution}

\begin{solution}{exo:g8:ohms-law:8}
\omterm{def:g7:short-circuits-safety:device}{Short circuit}: $R$ near zero, so $I = U/R$ explodes --- the
stampede is division by almost-nothing. Open \omterm{ex:g3:first-electric-circuit:switch}{switch}: $R$
beyond \omterm{def:g6:measurement-in-science:measuring}{measure}, so $I = U/R$ collapses to zero --- the total
blockade.
\end{solution}

\begin{solution}{exo:g8:ohms-law:9}
The \omterm{def:g8:ohms-law:resistor}{resistor} must take $9 - 2 = \qty{7}{V}$ at \qty{0.02}{A}:
$R = 7 \div 0.02 = 350\,\Omega$.
\end{solution}

\begin{solution}{exo:g8:ohms-law:10}
First point: $R = 1 \div 0.25 = 4\,\Omega$; second: $4 \div
0.5 = 8\,\Omega$ --- no single value fits. The story: the
growing current \omterm{def:g5:heat-and-insulation:heat}{heats} the filament, and hot metal resists
more; the portrait bends because the component changes as it
works.
\end{solution}

\begin{solution}{exo:g8:ohms-law:11}
One loop, one $I$: the two drops sum to the \omterm{def:g7:short-circuits-safety:battery}{battery's} push,
$6 = 10 I + 20 I = 30 I$, so $I = \qty{0.2}{A}$. Then $U_{10}
= 10 \times 0.2 = \qty{2}{V}$ and $U_{20} = 20 \times 0.2 =
\qty{4}{V}$ --- summing to \qty{6}{V}, as the loop law
demands.
\end{solution}

\begin{solution}{exo:g8:ohms-law:12}
Each \omterm{def:g5:series-parallel-circuits:parallel}{branch} spans \qty{6}{V}: $I_{10} = 6 \div 10 =
\qty{0.6}{A}$, $I_{20} = 6 \div 20 = \qty{0.3}{A}$; total
$\qty{0.9}{A}$. The stand-in: $R = 6 \div 0.9 \approx
6.7\,\Omega$ --- \emph{less} than either member. In
road-language: two roads between the same towns carry more
traffic than either alone; every added parallel road eases
the passage, so the team's opposition falls below its
weakest member's.
\end{solution}

\begin{solution}{pb:g8:ohms-law:1}
\textbf{1.} Supply, \omterm{def:g8:intensity-ammeter:ammeter}{ammeter} and component in one series loop
--- the tollbooth in the road; the \omterm{def:g8:voltage-voltmeter:voltmeter}{voltmeter} across the
component --- the surveyor aside.
\textbf{2.} $R = 4.5 \div 0.15 = 30\,\Omega$.
\textbf{3.} Five percent of $33$ is about $1.65$: acceptable
from $31.35$ to $34.65\,\Omega$. Sample one's $30\,\Omega$
falls outside: rejected.
\textbf{4.} $I = 4.5 \div 150 = \qty{0.03}{A} =
\qty{30.0}{mA}$.
\textbf{5.} Every row returns $R = U/I = 30\,\Omega$ ($1.5
\div 0.05$, $3.0 \div 0.10$, \dots): constant ratio,
straight-line portrait --- one value fits all: $30\,\Omega$.
\textbf{6.} The ratios run $5$, $6.7$, $8.3$, $10\,\Omega$:
climbing with every step --- \omterm{def:g8:ohms-law:resistance}{resistance} growing as the
current grows; no single $R$ exists.
\textbf{7.} As the current grows the filament \omterm{def:g5:heat-and-insulation:heat}{heats} toward
glowing, and hot metal resists more --- the portrait bends
upward. The bench's straight-line test thus admits exactly
the components whose $R$ stays put: it is, by construction, a
resistor-detector.
\textbf{8.} The \omterm{def:g3:temperature-thermometers:temperature}{steady-temperature} clause: Ohm's law promises
proportionality \emph{at steady \omterm{def:g3:temperature-thermometers:temperature}{temperature}} --- warm a
component mid-portrait and even an honest \omterm{def:g8:ohms-law:resistor}{resistor} drifts.
\textbf{9.} $R = 12 \div 0.03 = 400\,\Omega$.
\textbf{10.} $U = R I = 25 \times 9.2 = \qty{230}{V}$ ---
exactly the mains: the coil draws precisely what Ohm's law
allots it, and the \omterm{def:g7:short-circuits-safety:fuse}{fuse}'s margin is working as designed.
Nothing is ``only''; all is arithmetic.
\textbf{11.} The jumper is born to be a connecting wire ---
the level road, spending no push. Misused across a supply, its
near-zero \omterm{def:g8:ohms-law:resistance}{ohms} make it the perfect \omterm{def:g7:short-circuits-safety:device}{short circuit}: the villain
was always just a very good wire in the wrong place.
\textbf{12.} For example: ``The law: $U = R I$, push equals
opposition times march. Its portrait: a straight line through
the origin, steeper for stronger opposers. And only a
component at steady \omterm{def:g3:temperature-thermometers:temperature}{temperature} may sign it --- \omterm{def:g5:heat-and-insulation:heat}{heat} rewrites
\omterm{def:g8:ohms-law:resistance}{resistance}, and portraits taken feverish prove nothing.''
\end{solution}
